\documentclass[10pt,a4paper]{amsart}
\usepackage{amsmath,amssymb,amsthm,graphicx}
% Retained preamble for the Stacks Project chapter snapshot in this directory.
% The chapter's own preamble.tex is a copy of the shared Stacks Project preamble; the block
% below is sliced MECHANICALLY, line for line, from it. Nothing was added, removed or
% redefined; the wrapper class, the xr-hyper/hyperref wiring and the externaldocument list
% are omitted because the assembler supplies the document class.
%
% lines taken from preamble.tex: 3, 5, 9, 11, 12, 156, 157, 158, ... (51 lines total)
%
% For commutative diagrams we use Xy-pic
\usepackage[all]{xy}

% We use 2cell for 2-commutative diagrams.
\xyoption{2cell}
\UseAllTwocells

\newenvironment{reference}{\comment}{\endcomment}
\newenvironment{slogan}{\comment}{\endcomment}
\newenvironment{history}{\comment}{\endcomment}
\theoremstyle{plain}
\newtheorem{theorem}[subsection]{Theorem}
\newtheorem{proposition}[subsection]{Proposition}
\newtheorem{lemma}[subsection]{Lemma}
\theoremstyle{definition}
\newtheorem{definition}[subsection]{Definition}
\newtheorem{example}[subsection]{Example}
\newtheorem{exercise}[subsection]{Exercise}
\newtheorem{situation}[subsection]{Situation}
\theoremstyle{remark}
\newtheorem{remark}[subsection]{Remark}
\newtheorem{remarks}[subsection]{Remarks}
\numberwithin{equation}{subsection}
\def\lim{\mathop{\mathrm{lim}}\nolimits}
\def\colim{\mathop{\mathrm{colim}}\nolimits}
\def\Spec{\mathop{\mathrm{Spec}}}
\def\Hom{\mathop{\mathrm{Hom}}\nolimits}
\def\Ext{\mathop{\mathrm{Ext}}\nolimits}
\def\SheafHom{\mathop{\mathcal{H}\!\mathit{om}}\nolimits}
\def\SheafExt{\mathop{\mathcal{E}\!\mathit{xt}}\nolimits}
\def\Sch{\mathit{Sch}}
\def\Mor{\mathop{\mathrm{Mor}}\nolimits}
\def\Ob{\mathop{\mathrm{Ob}}\nolimits}
\def\Sh{\mathop{\mathit{Sh}}\nolimits}
\def\NL{\mathop{N\!L}\nolimits}
\def\CH{\mathop{\mathrm{CH}}\nolimits}
\def\proetale{{pro\text{-}\acute{e}tale}}
\def\etale{{\acute{e}tale}}
\def\QCoh{\mathit{QCoh}}
\def\Ker{\mathop{\mathrm{Ker}}}
\def\Im{\mathop{\mathrm{Im}}}
\def\Coker{\mathop{\mathrm{Coker}}}
\def\Coim{\mathop{\mathrm{Coim}}}
\DeclareMathSymbol{\boxtimes}{\mathbin}{AMSa}{"02}
\def\QCohstack{\mathcal{QC}\!\mathit{oh}}
\def\Cohstack{\mathcal{C}\!\mathit{oh}}
\def\Spacesstack{\mathcal{S}\!\mathit{paces}}
\def\Quotfunctor{\mathrm{Quot}}
\def\Hilbfunctor{\mathrm{Hilb}}
\def\Curvesstack{\mathcal{C}\!\mathit{urves}}
\def\Polarizedstack{\mathcal{P}\!\mathit{olarized}}
\def\Complexesstack{\mathcal{C}\!\mathit{omplexes}}
\def\Pic{\mathop{\mathrm{Pic}}\nolimits}
\def\Picardstack{\mathcal{P}\!\mathit{ic}}
\def\Picardfunctor{\mathrm{Pic}}
\def\Deformationcategory{\mathcal{D}\!\mathit{ef}}
\usepackage{cleveref}
\crefname{originalreference}{original reference}{original references}
\Crefname{originalreference}{Original reference}{Original references}
\makeatletter
\newlabel{stacks-morphisms-lemma-composition-smooth}{{075V}{source}{}{sourcepointer-1}{}}
\newlabel{stacks-morphisms-lemma-composition-smooth@cref}{{[originalreference][0][]075V}{[1][0][]source}{}{}{}}
\makeatother
\crefname{originalreference}{original reference}{original references}
\Crefname{originalreference}{Original reference}{Original references}
\makeatletter
\newlabel{stacks-morphisms-lemma-composition-flat}{{06PX}{source}{}{sourcepointer-2}{}}
\newlabel{stacks-morphisms-lemma-composition-flat@cref}{{[originalreference][0][]06PX}{[1][0][]source}{}{}{}}
\makeatother
\crefname{originallemma}{original lemma}{original lemmas}
\Crefname{originallemma}{Original lemma}{Original lemmas}
\makeatletter
\newlabel{lemma-lisse-etale}{{0788}{source}{}{sourcepointer-3}{}}
\newlabel{lemma-lisse-etale@cref}{{[originallemma][0][]0788}{[1][0][]source}{}{}{}}
\makeatother
\crefname{originalreference}{original reference}{original references}
\Crefname{originalreference}{Original reference}{Original references}
\makeatletter
\newlabel{sites-lemma-composition-cocontinuous}{{03L5}{source}{}{sourcepointer-4}{}}
\newlabel{sites-lemma-composition-cocontinuous@cref}{{[originalreference][0][]03L5}{[1][0][]source}{}{}{}}
\makeatother
\crefname{originalreference}{original reference}{original references}
\Crefname{originalreference}{Original reference}{Original references}
\makeatletter
\newlabel{sites-section-pullback}{{06UM}{source}{}{sourcepointer-5}{}}
\newlabel{sites-section-pullback@cref}{{[originalreference][0][]06UM}{[1][0][]source}{}{}{}}
\makeatother
\crefname{originalreference}{original reference}{original references}
\Crefname{originalreference}{Original reference}{Original references}
\makeatletter
\newlabel{stacks-sheaves-equation-pushforward}{{06W6}{source}{}{sourcepointer-6}{}}
\newlabel{stacks-sheaves-equation-pushforward@cref}{{[originalreference][0][]06W6}{[1][0][]source}{}{}{}}
\makeatother
\crefname{originalreference}{original reference}{original references}
\Crefname{originalreference}{Original reference}{Original references}
\makeatletter
\newlabel{stacks-sheaves-lemma-cohomology-on-subcategory}{{07AK}{source}{}{sourcepointer-7}{}}
\newlabel{stacks-sheaves-lemma-cohomology-on-subcategory@cref}{{[originalreference][0][]07AK}{[1][0][]source}{}{}{}}
\makeatother
\title{Functoriality of the lisse-etale and flat-fppf sites under smooth and flat morphisms}
\author{The Stacks Project Authors}
\date{stacks-project:stacks-cohomology.tex:local-snapshot-20261012; GNU FDL 1.2 or later}
\begin{document}
\maketitle
\noindent\textit{Editorial note.} All delivered mathematical text is the author's own, copied verbatim from the pinned Stacks Project chapter file stacks-cohomology.tex; no statement or proof step was rewritten, reordered or removed.\par\medskip
\noindent\textit{Editorial note.} The author omits the verification of the assumptions of the cohomology-on-subcategory lemma, noting smoothness (resp. flatness) of V -> Y is what verifies the second condition.\par\medskip
\noindent\textit{Allowed standard prerequisites (editorial).}\begin{itemize}
\item Pushforward formula on stacks and cohomology on subcategories. Sheaves on Stacks, Equation stacks-sheaves-equation-pushforward; Sheaves on Stacks, Lemma stacks-sheaves-lemma-cohomology-on-subcategory; Sites, Lemma sites-lemma-composition-cocontinuous. Diagram and equality steps.
\end{itemize}\medskip
\ifdefined\hypertarget\hypertarget{sourcepointer-1}{}\fi
\noindent\textit{Source reference.} Original reference 075V: sources/stacks-chow/tags.txt, source lines 9144--9144. Cross-chapter standard prerequisite cited by the author inside the counted span; it is used as a pointer and its own proof is not part of this item.\par\medskip
\ifdefined\hypertarget\hypertarget{sourcepointer-2}{}\fi
\noindent\textit{Source reference.} Original reference 06PX: sources/stacks-chow/tags.txt, source lines 8586--8586. Cross-chapter standard prerequisite cited by the author inside the counted span; it is used as a pointer and its own proof is not part of this item.\par\medskip
\ifdefined\hypertarget\hypertarget{sourcepointer-3}{}\fi
\noindent\textit{Source reference.} Original lemma 0788: sources/stacks-chow/tags.txt, source lines 9228--9228. Cross-chapter standard prerequisite cited by the author inside the counted span; it is used as a pointer and its own proof is not part of this item.\par\medskip
\ifdefined\hypertarget\hypertarget{sourcepointer-4}{}\fi
\noindent\textit{Source reference.} Original reference 03L5: sources/stacks-chow/tags.txt, source lines 4701--4701. Cross-chapter standard prerequisite cited by the author inside the counted span; it is used as a pointer and its own proof is not part of this item.\par\medskip
\ifdefined\hypertarget\hypertarget{sourcepointer-5}{}\fi
\noindent\textit{Source reference.} Original reference 06UM: sources/stacks-chow/tags.txt, source lines 8751--8751. Cross-chapter standard prerequisite cited by the author inside the counted span; it is used as a pointer and its own proof is not part of this item.\par\medskip
\ifdefined\hypertarget\hypertarget{sourcepointer-6}{}\fi
\noindent\textit{Source reference.} Original reference 06W6: sources/stacks-chow/tags.txt, source lines 8805--8805. Cross-chapter standard prerequisite cited by the author inside the counted span; it is used as a pointer and its own proof is not part of this item.\par\medskip
\ifdefined\hypertarget\hypertarget{sourcepointer-7}{}\fi
\noindent\textit{Source reference.} Original reference 07AK: sources/stacks-chow/tags.txt, source lines 9310--9310. Cross-chapter standard prerequisite cited by the author inside the counted span; it is used as a pointer and its own proof is not part of this item.\par\medskip
\begin{lemma}
\label{lemma-lisse-etale-functorial}
Let $f : \mathcal{X} \to \mathcal{Y}$ be a morphism of algebraic stacks.
\begin{enumerate}
\item If $f$ is smooth, then $f$ restricts to a continuous and cocontinuous
functor
$\mathcal{X}_{lisse,\etale} \to \mathcal{Y}_{lisse,\etale}$
which gives a morphism of ringed topoi fitting into the following
commutative diagram
$$
\xymatrix{
\Sh(\mathcal{X}_{lisse,\etale}) \ar[r]_{g'} \ar[d]_{f'} &
\Sh(\mathcal{X}_\etale) \ar[d]^f \\
\Sh(\mathcal{Y}_{lisse,\etale}) \ar[r]^g &
\Sh(\mathcal{Y}_\etale)
}
$$
We have $f'_*(g')^{-1} = g^{-1}f_*$ and $g'_!(f')^{-1} = f^{-1}g_!$.
\item If $f$ is flat, then $f$ restricts to a continuous and cocontinuous
functor
$\mathcal{X}_{flat,fppf} \to \mathcal{Y}_{flat,fppf}$
which gives a morphism of ringed topoi fitting into the following
commutative diagram
$$
\xymatrix{
\Sh(\mathcal{X}_{flat,fppf}) \ar[r]_{g'} \ar[d]_{f'} &
\Sh(\mathcal{X}_{fppf}) \ar[d]^f \\
\Sh(\mathcal{Y}_{flat,fppf}) \ar[r]^g &
\Sh(\mathcal{Y}_{fppf})
}
$$
We have $f'_*(g')^{-1} = g^{-1}f_*$ and $g'_!(f')^{-1} = f^{-1}g_!$.
\end{enumerate}
\end{lemma}

\begin{proof}
The initial statement comes from the fact that if $x \in \Ob(\mathcal{X})$
lies over a scheme $U$ such that $x : U \to \mathcal{X}$ is smooth
(resp.\ flat) and if $f$ is smooth (resp.\ flat) then
$f(x) : U \to \mathcal{Y}$ is smooth (resp.\ flat), see
Morphisms of Stacks, Lemmas \ref{stacks-morphisms-lemma-composition-smooth} and
\ref{stacks-morphisms-lemma-composition-flat}. The induced functor
$\mathcal{X}_{lisse,\etale} \to \mathcal{Y}_{lisse,\etale}$
(resp.\ $\mathcal{X}_{flat,fppf} \to \mathcal{Y}_{flat,fppf}$) is
continuous and cocontinuous by our definition of coverings in these
categories. Finally, the commutativity of the diagram is a consequence of the
fact that the horizontal morphisms are given by the inclusion functors (see
Lemma \ref{lemma-lisse-etale}) and
Sites, Lemma \ref{sites-lemma-composition-cocontinuous}.

\medskip\noindent
To show that $f'_*(g')^{-1} = g^{-1}f_*$ let $\mathcal{F}$ be a sheaf
on $\mathcal{X}_\etale$ (resp.\ $\mathcal{X}_{fppf}$).
There is a canonical pullback map
$$
g^{-1}f_*\mathcal{F} \longrightarrow f'_*(g')^{-1}\mathcal{F}
$$
see Sites, Section \ref{sites-section-pullback}.
We claim this map is an isomorphism.
To prove this pick an object $y$ of $\mathcal{Y}_{lisse,\etale}$
(resp.\ $\mathcal{Y}_{flat,fppf}$). Say $y$ lies over the scheme $V$
such that $y : V \to \mathcal{Y}$ is smooth (resp.\ flat). Since
$g^{-1}$ is the restriction we find that
$$
\left(g^{-1}f_*\mathcal{F}\right)(y) =
\Gamma(V \times_{y, \mathcal{Y}} \mathcal{X},\ \text{pr}^{-1}\mathcal{F})
$$
by Sheaves on Stacks, Equation (\ref{stacks-sheaves-equation-pushforward}).
Let
$(V \times_{y, \mathcal{Y}} \mathcal{X})' \subset
V \times_{y, \mathcal{Y}} \mathcal{X}$
be the full subcategory consisting of objects
$z : W \to V \times_{y, \mathcal{Y}} \mathcal{X}$ such that the induced
morphism $W \to \mathcal{X}$ is smooth (resp.\ flat). Denote
$$
\text{pr}' :
(V \times_{y, \mathcal{Y}} \mathcal{X})'
\longrightarrow
\mathcal{X}_{lisse,\etale}
\ (\text{resp. }\mathcal{X}_{flat,fppf})
$$
the restriction of the functor $\text{pr}$ used in the formula above.
Exactly the same argument that proves
Sheaves on Stacks, Equation (\ref{stacks-sheaves-equation-pushforward})
shows that for any sheaf $\mathcal{H}$ on $\mathcal{X}_{lisse,\etale}$
(resp.\ $\mathcal{X}_{flat,fppf}$) we have
\begin{equation}
\label{equation-pushforward-lisse-etale}
f'_*\mathcal{H}(y) =
\Gamma((V \times_{y, \mathcal{Y}} \mathcal{X})',
\ (\text{pr}')^{-1}\mathcal{H})
\end{equation}
Since $(g')^{-1}$ is restriction we see that
$$
\left(f'_*(g')^{-1}\mathcal{F}\right)(y) =
\Gamma((V \times_{y, \mathcal{Y}} \mathcal{X})',
\ \text{pr}^{-1}\mathcal{F}|_{(V \times_{y, \mathcal{Y}} \mathcal{X})'})
$$
By
Sheaves on Stacks, Lemma \ref{stacks-sheaves-lemma-cohomology-on-subcategory}
we see that
$$
\Gamma((V \times_{y, \mathcal{Y}} \mathcal{X})',
\ \text{pr}^{-1}\mathcal{F}|_{(V \times_{y, \mathcal{Y}} \mathcal{X})'})
=
\Gamma(V \times_{y, \mathcal{Y}} \mathcal{X},\ \text{pr}^{-1}\mathcal{F})
$$
are equal as desired; although we omit the verification of the assumptions
of the lemma we note that the fact that $V \to \mathcal{Y}$ is smooth
(resp.\ flat) is used to verify the second condition.

\medskip\noindent
Finally, the equality $g'_!(f')^{-1} = f^{-1}g_!$ follows formally from
the equality $f'_*(g')^{-1} = g^{-1}f_*$ by the adjointness of
$f^{-1}$ and $f_*$, the adjointness of $g_!$ and $g^{-1}$, and their
``primed'' versions.
\end{proof}

\end{document}
